\documentclass{butexFR}

%-------------------- Informations ----------------------
\titre{Solution aux besoins}
\UE{Analyse}
\sujet{SAÉ 3.02 - Labyrinthe}
\enseignant{Mme Boneva}
\eleves{Bredelle Victor \\ Marouze Antonin \\ Harlaut Romain \\ Zheng Angèl \\ Lavogiez Baptiste}

%-------------------- Début du document -----------------
\begin{document}

\fairepagedegarde
\fairemarges
\tabledematieres


%======================= Contenu ========================

Qst à poser :

cas d'utilisation : unicité par fiche, regroupement ?
Diagramme de classe : controller ?


\section{Introduction}

Le présent rapport détaille la composante \texttt{Analyse} relative au projet \textbf{SAÉ S3.02 - Labyrinthe}. Il aborde l'organisation interne de notre équipe ainsi que notre approche de modélisation des exigences du client, structurée comme suit :

\begin{enumerate}
    \item Diagrammes de cas d'utilisation
    \item Fiches descriptives
    \item Prototype d'interface
    \item Diagramme de classe
\end{enumerate}

Il est entendu que chaque section est itérative.

\section{Équipe et organisation du travail}
\textbf{Groupe G3.3}

\begin{itemize}
    \item Bredelle Victor
    \item Marouze Antonin
    \item Harlaut Romain
    \item Zheng Angèl
    \item Lavogiez Baptiste
\end{itemize}

\tsec{Victor}

Rédaction de trois fiches descriptives avec Baptiste :
\begin{enumerate}
    \item Charger profil
    \item Afficher progression
    \item Lancer mode progression
\end{enumerate}

\tsec{Antonin}

Réalisation de la maquette \textit{Figma} avec Angèl / Romain

\tsec{Romain}

Réalisation de la fiche descriptive "Avancer d'une case", aide à la maquette \textit{Figma}

\tsec{Angèl}

Réalisation de la maquette \textit{Figma} avec Antonin / Romain

\tsec{Baptiste}

Rédaction de trois fiches descriptives avec Victor :
\begin{enumerate}
    \item Charger profil
    \item Afficher progression
    \item Lancer mode progression
\end{enumerate}

Rédaction du présent rapport

\tsecnonum{Le groupe}

Outre les responsabilités individuelles, nous avons collectivement débattu de la manière de modéliser chaque exigence en l'examinant sous divers angles.


\info{\textbf{L'ensemble des contributeurs a collaboré à l'élaboration du diagramme de classes.}}

\section{Cas d'utilisation}

\subsection{Diagrammes de cas d'utilisation}

screen (Review en tp avec autres)

\subsection{Fiches descriptives}

\subsubsection{Cas d’utilisation: Charger profil}
\textbf{Système:} Labyrinthe\\
\textbf{Acteur principal:} Utilisateur\\
\textbf{Déclencheur:} /\\
\textbf{Autres acteurs:} /\\[0.25em]
\textbf{Préconditions:} L'utilisateur est dans le menu de saisie du nom (``Lancer le mode progression'')\\
\textbf{Garanties en cas de succès:} Le système charge le profil associé au nom et l'utilisateur reprend sa progression\\
\textbf{Garanties minimales:} /\\[0.25em]
\textbf{Scénario nominal:}
\begin{enumerate}
    \item L'utilisateur a rentré un nom déjà existant dans la base de données lors du \textbf{Cas d'utilisation} ``Lancer le mode progression'' (suite de ce cas, ayant déjà vérifié l'occurence)
    \item Le système affiche un menu où l'utilisateur peut choisir de ``Reprendre la partie existante'' ou ``Écraser la partie existante''.
    \item L'utilisateur sélectionne ``Reprendre la partie''
    \item Le système charge le profil associé au nom depuis la base de données et l'utilisateur reprend sa progression
\end{enumerate}
\textbf{Scénario alternatifs:} /

\newpage

\subsubsection{Cas d’utilisation: Afficher progression}
\textbf{Système:} Labyrinthe\\
\textbf{Acteur principal:} Utilisateur\\
\textbf{Déclencheur:} /\\
\textbf{Autres acteurs:} /\\[0.25em]
\textbf{Préconditions:} Avoir un profil, être dans le mode progression\\
\textbf{Garanties en cas de succès:} Le système affiche la progression associée au joueur (score, caractéristiques...)\\
\textbf{Garanties minimales:} /\\[0.25em]
\textbf{Scénario nominal:}
\begin{enumerate}
    \item L'utilisateur sélectionne la fonctionnalité ``accéder au menu du mode progression''
    \item Le système récupère les données du joueur dans la base et calcule le score et caractéristiques afin d'afficher la progression
\end{enumerate}
\textbf{Scénario alternatifs:} /

\newpage

\subsubsection{Cas d’utilisation: Lancer mode progression}
\textbf{Système:} Labyrinthe\\
\textbf{Acteur principal:} Utilisateur\\
\textbf{Déclencheur:} /\\
\textbf{Autres acteurs:} /\\[0.25em]
\textbf{Préconditions:} / (avoir lancé le jeu)\\
\textbf{Garanties en cas de succès:} L'utilisateur arrive sur le mode progression\\
\textbf{Garanties minimales:} /\\[0.25em]
\textbf{Scénario nominal:}
\begin{enumerate}
    \item L'utilisateur sélectionne la fonctionnalité ``Lancer mode progression''.
    \item Le système affiche un menu où l'utilisateur peut rentrer son nom
    \item L'utilisateur rentre son nom
    \item Le système vérifie que le nom est valide, sauvegarde son profil dans la base de données et lance le mode progression
\end{enumerate}
\textbf{Scénario alternatifs:}

\emph{A: Le système détecte une erreur lors de la vérification}\\
\hspace*{1em}4(A): Le système affiche une notification (de quelles mauvaises informations il s'agit), puis revient en étape 2 du scénario principal

\emph{B: Le système détecte une occurence du nom dans la base de données (charger partie)}\\
\hspace*{1em}4(B): Le système démarre le \textbf{Cas d'utilisation} ``Charger profil''

\newpage

\subsubsection{Cas d’utilisation: Déplacer joueur}
\textbf{Système:} Labyrinthe\\
\textbf{Acteur principal:} Utilisateur\\
\textbf{Déclencheur:} /\\
\textbf{Autres acteurs:} /\\[0.25em]
\textbf{Préconditions:} Être dans le labyrinthe\\
\textbf{Garanties en cas de succès:} Le système a déplacé le joueur à la case souhaitée\\
\textbf{Garanties minimales:} /\\[0.25em]
\textbf{Scénario nominal:}
\begin{enumerate}
    \item Le joueur appuie sur une flèche directionnelle
    \item Le système vérifie que la direction est accessible et déplace le joueur sur la case associée
\end{enumerate}
\textbf{Scénario alternatifs:}

\emph{A: Le joueur arrive sur la sortie}\\
\hspace*{1em}2(A): Le système affiche une notification annonçant que la sortie a été atteinte et redirige vers le menu principal.

\emph{B: Le joueur veut se déplacer dans une case inaccessible}\\
\hspace*{1em}2(B): Le système affiche la position actuelle et pas la nouvelle position et revient en étape 1 du scénario principal

\newpage

\subsection{Prototypes d'interface}

\begin{figure}[h!]
    \centering
    \includegraphics[width=1\textwidth]{figures/cas_dutilisation/prototypes_dinterface/default/fig_1.png}
    \caption{Prototype principal}
    \label{fig:cas_dutilisation_prototypes_dinterface_1}
\end{figure}

Faire photos successives

L'interface est pensée pour être simple, intuitive, et sans ambiguité. La priorité demeure, comme toujours, le \textbf{strict respect des attentes du client}.

\section{Diagramme de classe}

Figure

Le diagramme de classe a été généré à l'aide de \textit{PlantUML} dans un format textuel, assurant ainsi une portabilité et une facilité de modification accrues, quelle que soit la localisation.

La conception de ce diagramme a nécessité un investissement conséquent, considérant que, par nature, un diagramme de classe requiert une \textbf{réflexion soutenue et l'intégration de multiples perspectives}. 

Des retours ont été sollicités auprès des professeurs spécialisés en analyse, ainsi qu'auprès de ceux enseignant la qualité de développement.

Par conséquent, nous estimons avoir abouti à une solution hautement optimisée, caractérisée par l'absence de redondance et le respect du \textbf{principe de responsabilité unique} pour chaque classe. 

\merci

\end{document}




